% !TITLE 示儿
% !AUTHOR 陆游
% !DYNASTY 两宋
% !GENRE 七言绝句
% !ORDER 18
% !EDITION 《陆游集》，中华书局（1976年版为候选）
% !EDITION_STATUS pending

\begin{poem}
死去元知\textsuperscript{1}万事空，但悲\textsuperscript{2}不见九州同\textsuperscript{3}。
王师北定中原日\textsuperscript{4}，家祭无忘告乃翁\textsuperscript{5}。
\end{poem}

\begin{annotations}
\notetext{1}{元知：本来就知道。元，同“原”。人死之后万事成空，这个道理谁都明白。}
\notetext{2}{但悲：只悲伤。但，只、只是。明白万事皆空，却仍有一件事放不下。}
\notetext{3}{九州同：全国统一。九州是古人对中国大地的划分，后来代指中国；“同”指收复中原、天下一统。}
\notetext{4}{王师北定中原日：等到朝廷的军队北上收复中原的那一天。王师，指南宋的军队；中原，指被金人占领的黄河中下游一带。}
\notetext{5}{家祭无忘告乃翁：家里祭祖的时候，别忘了告诉我。无，同“毋”，不要；乃翁，你们的父亲。这是陆游八十五岁临终前写给儿子们的绝笔诗——他没能等到那一天，只好把这件事托付给身后。}
\end{annotations}

\begin{appreciation}
这是陆游的绝笔，写于八十五岁临终之前，是中国诗史上最著名的遗嘱之一。

“死去元知万事空”——本来就知道，人一死什么都成了空。“元”同“原”，“元知”两字是说这件事他早就想明白了，开篇一句便把话说到了底：既然万事皆空，那还有什么可说的？偏偏下一句就转了过来。

“但悲不见九州同”——只悲伤一件事，没能看见山河一统。九州是古人对中国大地的划分，后来代指中国；“同”指收复中原、天下归一。一个“但”字是全诗的支点，把前面那句“万事空”整个掀翻：万事都能空，唯独这一件空不了。

“王师北定中原日，家祭无忘告乃翁”——等到朝廷的军队北上收复中原的那一天，家里祭祖的时候，别忘了告诉我。王师是南宋的军队，中原是被金人占去的黄河中下游一带，乃翁就是“你们的父亲”。他把一桩身后的大事郑重地托付出去，托付给儿子，也托付给一个自己看不见的日子。

最动人的是那个“告”字。他明知万事皆空，明知自己听不见，还是要儿子们告诉他。这不是相信鬼神，是把遗愿说成一个约定——哪怕这边无人接应，话也要说出去，一代一代说下去。

陆游一生留诗九千余首，写恢复中原的不下千首。八十五岁，油尽灯枯，最后一点力气仍旧用在北伐这件事上。后来文天祥在狱中写“人生自古谁无死，留取丹心照汗青”，那是对自己说的，说的是名节；陆游这一句是对儿子说的，说的是国家的将来。一个把生死看淡了，一个把生死之后的事还紧紧攥在手里。

诗写完不久，他就去世了。此后的中原始终没有回到宋人手里，这句嘱托便一直悬着，没有等到回音。
\end{appreciation}
